\documentclass[../../main.tex]{subfiles}
\graphicspath{{../../image/}{../../figure/}}
\begin{document}

\section{Cauchy 问题：弱解}\label{section:ch04-弱解}

鉴于上一节所展示的古典解崩溃的例子--以及全书将遇到的更多例子--考虑守恒律组\eqref{align:ch04-守恒律组}的弱解变得势在必行。弱解的自然概念应与存在性理论相配合而确定。标量守恒律的弱解存在性问题，在任意多个空间变量下（见第 VI 章），以及至少在单个空间变量的方程组下（见第 XIII--XVII 章），已被以确定方式解决；然而，对多个空间变量的方程组，它仍完全未决。在缺乏确定性存在性理论的情况下，为引入若干相关概念而不对通量函数 $G$ 施加技术性增长条件，我们在此把\eqref{align:ch04-守恒律组}的\emph{弱解}定义为：局部有界、可测的函数 $U$，定义在 $\mathbb{R}^{m}\times[0,T)$ 上、取值于 $\mathcal{O}$，并在分布意义下满足\eqref{align:ch04-守恒律组}。

回顾\cref{chapter:ch01-平衡律}，我们把\eqref{align:ch04-守恒律组}的任一弱解 $U$ 规范化，使得映射 $t\mapsto U(\cdot,t)$ 在 $L^{\infty}(\mathbb{R}^{m})$ weak$^{*}$ 下于 $[0,T)$ 连续。规范化的弱解 $U$ 若还在 $\mathbb{R}^{m}$ 上几乎处处满足\eqref{align:ch04-Cauchy初值}，则它求解 Cauchy 问题\eqref{align:ch04-守恒律组}、\eqref{align:ch04-Cauchy初值}。\cref{chapter:ch01-平衡律}还蕴含

\begin{align}
\int_{\tau}^{T}\int_{\mathbb{R}^{m}}\left[\partial_{t}\Phi U+\sum_{\alpha=1}^{m}\partial_{\alpha}\Phi G_{\alpha}(U)\right]\,\mathrm{d}x\mathrm{d}t+\int_{\mathbb{R}^{m}}\Phi(x,\tau)U(x,\tau)\,\mathrm{d}x=0,
\label{align:ch04-弱解tau}
\end{align}
它对每一在 $\mathbb{R}^{m}\times[0,T)$ 内紧支集、取值于 $\mathbb{M}^{1\times n}$ 的 Lipschitz 检验函数 $\Phi(x,t)$ 以及任一 $\tau\in[0,T)$ 成立。特别地，
\begin{align}
\int_{0}^{T}\int_{\mathbb{R}^{m}}\left[\partial_{t}\Phi U+\sum_{\alpha=1}^{m}\partial_{\alpha}\Phi G_{\alpha}(U)\right]\,\mathrm{d}x\mathrm{d}t+\int_{\mathbb{R}^{m}}\Phi(x,0)U_{0}(x)\,\mathrm{d}x=0,
\label{align:ch04-弱解0}
\end{align}
它可作为\eqref{align:ch04-守恒律组}、\eqref{align:ch04-Cauchy初值}弱解的等价定义。$t\mapsto U(\cdot,t)$ 的连续性还导出所期望的半群性质：若 $U(x,t)$ 是\eqref{align:ch04-守恒律组}在 $[0,T)$ 上以 $U(x,0)$ 为初值的弱解，则对任一 $\tau\in[0,T)$，函数 $U_{\tau}(x,t)=U(x,t+\tau)$ 是以 $U_{\tau}(x,0)=U(x,\tau)$ 为初值的\eqref{align:ch04-守恒律组}的弱解。

一类重要的弱解是 $U$ 在 $\mathbb{R}^{m}\times[0,T)$ 上局部有界变差的那些。这类解在测度意义下满足方程组\eqref{align:ch04-守恒律组}，并把初始条件\eqref{align:ch04-Cauchy初值}作为 $U$ 在 $t=0$ 处的迹。在近似跳跃间断点集上，$BV$ 解满足 Rankine--Hugoniot 跳跃条件\cref{chapter:ch03-双曲平衡律方程组}，此处取形式

\begin{align}
-s[U_{+}-U_{-}]+[G(U_{+})-G(U_{-})]\nu=0.
\label{align:ch04-RH}
\end{align}

由于方程组是散度形式的，存在一个从空间变量到时间变量的正则性传递机制，可在 $BV$ 解的框架内加以说明：

\begin{theorem}\label{theorem:ch04-BV时间Lipschitz}
设 $U$ 是\eqref{align:ch04-守恒律组}在 $[0,T)$ 上的有界弱解，使得对任一固定 $t\in[0,T)$，$U(\cdot,t)\in BV(\mathbb{R}^{m})$，且 $\operatorname{TV}_{\mathbb{R}^{m}}U(\cdot,t)\leqslant V$（对所有 $t\in[0,T)$）。则 $t\mapsto U(\cdot,t)$ 在 $L^{1}(\mathbb{R}^{m})$ 中于 $[0,T)$ 上 Lipschitz 连续，
\begin{align}
\|U(\cdot,\sigma)-U(\cdot,\tau)\|_{L^{1}(\mathbb{R}^{m})}\leqslant aV|\sigma-\tau|,\qquad 0\leqslant\tau<\sigma<T,
\label{align:ch04-BVLipschitz}
\end{align}
其中 $a$ 仅依赖于 $G$ 的 Lipschitz 常数。特别地，$U$ 在 $\mathbb{R}^{m}\times[0,T)$ 上属于 $BV_{\mathrm{loc}}$。
\end{theorem}

\begin{proof}
固定 $0\leqslant\tau<\sigma<T$ 以及任一 $\Psi\in C_{0}^{\infty}(\mathbb{R}^{m};\mathbb{M}^{1\times n})$，且当 $x\in\mathbb{R}^{m}$ 时 $|\Psi(x)|\leqslant 1$。取检验函数 $\Phi(x,t)=f(t)\Psi(x)$，其中 $f\in C_{0}^{0}[0,T)$，且当 $t\in[0,\sigma]$ 时 $f(t)=1$。先对 $\tau=\sigma$、再对 $\tau=\tau$ 写出\eqref{align:ch04-弱解tau}，相减得
\begin{align}
\int_{\mathbb{R}^{m}}\Psi(x)[U(x,\sigma)-U(x,\tau)]\,\mathrm{d}x=\int_{\tau}^{\sigma}\int_{\mathbb{R}^{m}}\sum_{\alpha=1}^{m}\partial_{\alpha}\Psi(x)G_{\alpha}(U(x,t))\,\mathrm{d}x\mathrm{d}t.
\label{align:ch04-时间差}
\end{align}
右端空间积分由 $G(U(\cdot,t))$ 的全变差控制，而后者又以 $aV$ 为界。对一切满足 $|\Psi(x)|\leqslant 1$ 的 $\Psi$ 取上确界，即得\eqref{align:ch04-BVLipschitz}。

\cref{chapter:ch01-平衡律}与\eqref{align:ch04-BVLipschitz}合起来蕴含 $U$ 在 $\mathbb{R}^{m}\times[0,T)$ 上属于 $BV_{\mathrm{loc}}$。证明完成。

\end{proof}

\end{document}
